\begin{principle}[Three coefficient strata]
\label{prin:zeckendorf-three-coefficient-strata}
The three Window-six coefficients
$$
  \frac89,\qquad \frac23,\qquad \frac49
$$
are not competing fits for the same low-energy scalar. They occupy
different physical strata of the same finite readout geometry. The
phase-reservoir coefficient
$$
  q_{\mathrm{static}}^{\mathrm{imp}}=\frac89
$$
is the static impedance coefficient selected by the right-boundary
flux-reservoir readout. The image coefficient
$$
  q_{\beta}^{Z_6}=\frac23
$$
is the six-element $Z_6\to U(1)$ phase-image coefficient when that image
is read as a unit electromagnetic beta-flow activation. The one-way
coefficient
$$
  q_{\mathrm{curv}}^{\mathrm{one}}=\frac49
$$
is the irreducible curvature-residual or one-way fold-mismatch
coefficient.
\end{principle}

\begin{theorem}[Window-six readout diamond]
\label{thm:zeckendorf-window-six-readout-diamond}
The Window-six readout coefficients form the three-layer diamond
$$
  q_{\mathrm{curv}}
  =
  \frac49
  \quad\longrightarrow\quad
  q_{\beta}
  =
  \frac69
  =
  \frac23
  \quad\longrightarrow\quad
  q_{\mathrm{imp}}
  =
  \frac89.
$$
The common step is
$$
  \frac29
  =
  \frac{|P_6^{\mathrm{geo}}|}{N_{\geq4}(6)}.
$$
Thus the layers differ by activation of one basic $Z_2$ parity unit
against the deepest residual-tail denominator.
\end{theorem}
\begin{proof}
The curvature-residual layer quotients the locally absorbable diagonal
parity
$$
  P_6^{\mathrm{geo}}\cong\mathbb{Z}/2\mathbb{Z}
$$
from the boundary reservoir, so its visible residual group has order
$4$ and coefficient $4/N_{\geq4}(6)=4/9$. The $U(1)$ phase-image layer
uses the six-element image of the
$C_\partial\times T\to\mathbb{Z}/6\mathbb{Z}\to U(1)$ interface, so its
coefficient is $6/9=2/3$. The static impedance layer reads the full
boundary phase reservoir $C_\partial\cong(\mathbb{Z}/2\mathbb{Z})^3$,
so its coefficient is $8/9$. Since the orders are $4,6,8$, the
successive normalized differences are both $2/9$.
\end{proof}

\begin{theorem}[Static impedance selects $8/9$ while beta flow selects $2/3$]
\label{thm:zeckendorf-static-impedance-beta-flow-coefficient-split}
The low-energy denominator
$$
  D_{6,\alpha}^{\star}
  =
  47+\varphi^{-7}-\frac12\varphi^{-17}
  +\frac89\varphi^{-27}
$$
uses $8/9$ because the infrared electromagnetic scalar is an impedance
readout of the phase reservoir. The abelian unit-charge shell flow uses
$2/3$ because the running layer reads the $Z_6\to U(1)$ phase image.
\end{theorem}
\begin{proof}
By
\autoref{thm:zeckendorf-fixed-electromagnetic-readout-selection}, the
fixed infrared electromagnetic readout selects the phase-reservoir
coefficient $q_2=8/9$, not the image coefficient or the curvature
quotient. By
\autoref{prin:zeckendorf-right-block-flux-reservoir-coefficient}, the
same $8/9$ is the normalized right-block escape flux
$P(\mathrm{leave}\ U_R\mid U_R)=32/36$. This is the static impedance
layer.

By contrast, the $Z_6\to U(1)$ image has six phase images over the
Window-six depth denominator $N_{\geq4}(6)=9$, hence coefficient $6/9$.
In the shell-flow dictionary of
\autoref{thm:zeckendorf-renormalization-dictionary-golden-shell-flow},
an abelian unit Dirac charge contributes
$$
  \frac{d\gamma_\varphi}{d\log Q}
  =
  \frac{2}{3}
  \frac{\alpha}{\pi\log\varphi}.
$$
Thus the $2/3$ image coefficient is naturally assigned to the beta-flow
layer rather than to the static denominator.
\end{proof}

\begin{definition}[Window-six beta-shell coefficient]
\label{def:zeckendorf-window-six-beta-shell-coefficient}
The unit $Z_6$ beta-shell coefficient is
$$
  q_{\beta}^{Z_6}:=\frac23.
$$
For a scale interval with active electromagnetic modes, write
$$
  \mathcal{B}(Q)
  =
  q_{\beta}^{Z_6}
  \sum_f N_c(f)q_f^2\Theta_f(Q)
  +
  \mathcal{B}_{\mathrm{bos}}(Q)
  +
  \mathcal{B}_{\mathrm{had}}(Q),
$$
where $\Theta_f(Q)$ records threshold activation, $N_c(f)$ is the color
multiplicity, and $q_f$ is the electric charge in units of the
positron charge. The corresponding golden shell flow is
$$
  \frac{d\gamma_{\mathrm{eff}}}{d\log Q}
  =
  \frac{\alpha(Q)}{\pi\log\varphi}\,
  \mathcal{B}(Q).
$$
\end{definition}

\begin{corollary}[Coefficient layer separation]
\label{cor:zeckendorf-coefficient-layer-separation}
The low-energy static layer, the beta-shell layer, and the Green
susceptibility layer are separated as follows:
$$
  \frac89
  \quad\text{sets the static flux-reservoir boundary condition,}
$$
$$
  \frac23
  \quad\text{sets the unit }Z_6\text{ phase-image beta coefficient,}
$$
and
$$
  571
  \quad\text{certifies the internal Green susceptibility denominator.}
$$
The curvature coefficient $4/9$ remains the one-way mismatch or
curvature-residual coefficient. It does not replace either the static
impedance coefficient or the unit beta-flow coefficient.
\end{corollary}
